\documentclass[11pt]{article}
\usepackage{amsmath,amssymb,amsthm,hyperref}

\title{QPI Sieve: A Prime Classification Algorithm with Receipts}
\author{Christopher Michael LeBlanc Sr.}
\date{\today}

\newtheorem{theorem}{Theorem}
\newtheorem{definition}{Definition}
\newtheorem{remark}{Remark}

\begin{document}
\maketitle

\begin{abstract}
This paper documents a prime classification framework (QPI) and provides
reproducible computational receipts verifying the Twin Prime Midpoint Law
up to a finite bound $N$. All infinite-claim statements are explicitly
labeled as conjectures and are not asserted as theorems.
\end{abstract}

\section{Definitions}
\begin{definition}[Yellow]
An integer $m$ is \emph{Yellow} iff $3 \mid m$ and the last digit of $m$ is in $\{0,2,8\}$.
\end{definition}

\begin{definition}[Red]
An integer $m$ is \emph{Red} iff $m$ is Yellow and both $m-1$ and $m+1$ are prime.
\end{definition}

\begin{definition}[Twin prime]
A \emph{twin prime pair} is $(p,p+2)$ where $p$ and $p+2$ are prime.
\end{definition}

\section{Twin Prime Midpoint Law (Finite, Checkable)}
\begin{theorem}[Twin Prime Midpoint Law]
For integers $p>5$, $p$ and $p+2$ are prime if and only if $m=p+1$ is Red.
\end{theorem}

\begin{remark}
This theorem is an equivalence between:
(i) twin primality and (ii) a midpoint predicate consisting of a modular filter (Yellow)
plus neighbor primality. It does not imply an infinite number of twins.
\end{remark}

\section{Receipts}
The repository includes JSON receipts and SHA-256 hashes for runs at:
$N=10^6,10^7,10^8$. Each receipt records:
\begin{itemize}
\item \# twin pairs (all)
\item \# twin pairs with $p\ge 7$
\item \# red cells
\item the diff $\#\text{red}-\#\text{twin}(p\ge 7)$
\end{itemize}
A passing receipt has diff $=0$.
See \texttt{results/checksums.txt}.

\section{Formalization}
Lean 4 skeletons are provided in \texttt{lean/qpi\_theorems\_clean.lean}
with placeholders marked \texttt{sorry}. The infinitude claim is a \texttt{conjecture}.

\bibliographystyle{plain}
\bibliography{refs}
\end{document}
